\begin{table}[!htbp]\centering\scriptsize\setlength{\tabcolsep}{4pt}
\caption{Hausse annuelle des principaux perdants selon la durée de la transition (€ par an)}\label{tab:ircsg_trans}
\begin{tabular}{@{}lrrrr@{}}
\toprule
& Hausse finale & \multicolumn{3}{c}{Hausse chaque année, transition sur} \\ \cmidrule(l){3-5}
& & 3 ans & 5 ans & 10 ans \\
\midrule
Retraités, médiane & +1\,392 & 464 & 278 & 139 \\
Retraités, 90\textsuperscript{e} centile & +3\,423 & 1\,141 & 685 & 342 \\
Couples de retraités, 90\textsuperscript{e} centile & +5\,659 & 1\,886 & 1\,132 & 566 \\
Retraités, 99\textsuperscript{e} centile & +7\,423 & 2\,474 & 1\,485 & 742 \\
Tous les foyers, 90\textsuperscript{e} centile & +2\,245 & 748 & 449 & 224 \\
Retraité seul, pension de 16 000 € & +1\,050 & 350 & 210 & 105 \\
Retraité seul, pension de 20 000 € & +1\,390 & 463 & 278 & 139 \\
Retraité seul, pension de 25 000 € & +1\,679 & 560 & 336 & 168 \\
Couple de retraités, 45 000 € de pensions (60/40) & +1\,834 & 611 & 367 & 183 \\
Couple sans enfant à un seul revenu, 80 000 € & +3\,715 & 1\,238 & 743 & 372 \\
\bottomrule
\end{tabular}
\par\vspace{3pt}\parbox{\linewidth}{\raggedright\scriptsize\color{encre}
Convergence linéaire : chaque année, chaque foyer paie une part croissante du nouvel impôt et décroissante de l'ancien, ce qui tient la recette chaque année. Centiles : foyers représentatifs du modèle (la dispersion réelle des hausses est plus forte ; le 99\textsuperscript{e} centile en est un minorant). Cas-types : tableau~\ref{tab:ircsg_cas}.}
\end{table}
